\section*{Ukázka:}
\subsection*{I.}
„A~vaše jméno?“

„Nejmenuji se; jsem nikdo. Jsem hůř než nikdo, jsem krysař.“

Muž, který takto mluvil, stál se vzpřímenou hlavou před vraty domu, v~nichž se do soumraku bělala ženská postava. Hleděl na ni svýma temnýma, pátravýma očima. Byl vysoký a~štíhlý, štíhlejší ještě ve svém přiléhavém sametovém kabátci a~v~úzkých nohavicích. Jeho ruce byly drobné a~jemné jako ruce paní. Neměl při sobě ani zbraně, ani hole, ač se zdálo, že přichází z~daleka, po cestách, jež vždy jisty nebyly. Svíral zato něco dlouhého a~ozdobného, co budilo zvědavost ženy, s~kterou hovořil. Byla to píšťala cizokrajné práce, jaké dosud neviděla.

„Krysař,“ zasmála se žena ve dveřích. „Přicházíte včas do Hammeln. Není tu krysaře; zato krys je tu mnoho. Vysvětlete mi, krysaři, odkud se berou krysy? Nebývalo jich dříve, jak nám říkají. Je pravda,“ končila s úsměvem, „staří lidé míní, že svět je stále horší.“

Krysař pokrčil rameny.

„Odkud jsou, nevím. Jsou však v~každém z~vašich domů. Hlodají bez ustání; hlodají dole ve sklepích, hlodají tam, kde jich nevidíme. Stanou se však dotěrnými a~stoupají pak výše. Připravujete hostinu, svatbu, křtiny, což já vím? Představte si, že se náhle při hostině zjeví krysy s~dlouhými boltci a~dlouhými kníry. Tento zjev kazí chuť, chápete přece.“

„Ano,“ smála se žena ve dveřích. „Při svatbě Kateřiny objevila se pořádná krysa. Ženich byl bledý jako stěna a~Kateřina padla do mdlob. Lidé nesnesou nic tak málo jako to, co jim kazí chuť; odhodlají se pak zavolat krysaře.“

„Připravujete svatbu nebo křtiny?“ otázal se krysař náhle, bez přechodu.

Žena ve dveřích se zasmála hlasitě.

„Jste cizinec, je zřejmo, že jste cizinec. Nejsem vdána, krysaři.“ Krysař se uklonil.

„To nevadí… nijak to nevadí. Nuže, povolají krysaře. Krysař píská a~píská, až vyvede všechnu havěť z~jejích skrýší. Jde za ním jako omámená. Zavede ji do řeky, do Rýna, Dunaje, Havoly, Vesery. A~dům je prostý krys.“

Krysař opět se uklonil a jeho hlas zachvěl se nějakou elegií. Žena mlčela, pohrávajíc si větvičkou jasmínu.

„Ale je-li tomu tak, nepomyslí nikdo na krysaře. Krysař, cizinko, je muž, který nezůstává, ale jde. Lidé vidí ho rádi přicházet. Ale odcházeti vidí ho mnohem raději.“

„Opravdu?“ řekla pouze. Znělo to jako povzbuzení a~snad to ani povzbuzení nebylo. Ale krysař chápal to tak. Bledé jeho tváře se zbarvily; byla by to mohla pozorovat, nebýti temna.

„Cítím to, cizinko. Lidé nemilují krysaře; bojí se ho pouze.“ Děvče zasmálo se opět.

„A~čím to, že jdou krysy tak slepě za vámi, krysaři?“

Krysař ukázal na píšťalu, která oživla tímto pohybem.

„Je to zvláštní nástroj,“ řekl.

Pohlédla zvědavě na mluvčího i~na píšťalu. Dotkla se jí lehce.

„Píšťala,“ řekla pak pohrdavě. „Pěkná píšťala. Ale nic více než píšťala.“

„Krysy mají dobrý sluch a~má píšťala dobrý zvuk.“

Oči krysařovy zaplály podivným ohněm. Děvče ve dveřích maně ucouvlo. Větvička jasmínu v~její ruce se chvěla.

„Mám zvláštní nadání vyhánět krysy,“ pokračoval krysař. „Pískám občas na ni velmi smutné písně, písně všech krajů, kterými jsem prošel. A~já prošel mnoha zeměmi; slunnými i~ponurými, pláněmi a horami. Má píšťala zní tak tiše. Krysy slyší a~jdou. Není takového krysaře mimo mne. Povím vám něco, cizinko, která se umíte tak zvonivě smát. Nikdy jsem nepískal plným dechem; byl to vždy tlumený zvuk. Kdybych pískal plně, ne pouze krysy šly by za mnou.“

Krysař domluvil. Oči uhasly a svěsil bezděčně ruce s píšťalou.

\cite[4-5]{krysař}
